
\documentclass{article}
\usepackage{colortbl}
\usepackage{makecell}
\usepackage{multirow}
\usepackage{supertabular}

\begin{document}

\newcounter{utterance}

\centering \large Interaction Transcript for game `referencegame', experiment `random\_grids', episode 16 with reflexion{-}Qwen3.5{-}9B.
\vspace{24pt}

{ \footnotesize  \setcounter{utterance}{1}
\setlength{\tabcolsep}{0pt}
\begin{supertabular}{c@{$\;$}|p{.15\linewidth}@{}p{.15\linewidth}p{.15\linewidth}p{.15\linewidth}p{.15\linewidth}p{.15\linewidth}}
    \# & $\;$A & \multicolumn{4}{c}{Game Master} & $\;\:$B\\
    \hline

    \theutterance \stepcounter{utterance}  
    & & \multicolumn{4}{p{0.6\linewidth}}{
        \cellcolor[rgb]{0.9,0.9,0.9}{
            \makecell[{{p{\linewidth}}}]{
                \texttt{\tiny{[P1$\langle$GM]}}
                \texttt{You have attempted this same task before and did not succeed. These are the plans you wrote after those attempts.} \\
\texttt{Trial \#1: New plan: My previous expression failed because I misidentified the filled cells in the third row; I stated positions 1 and 4 were filled, but the target grid actually has filled cells at positions 1 and 4 in the third row? Wait, let me re{-}examine the target grid carefully.} \\
\texttt{Target Row 3: `X □ □ X □` {-}> Filled at 1 and 4. My description said 1 and 4. That seems correct.} \\
\texttt{Let's re{-}read the target grid vs my description.} \\
\texttt{Target:} \\
\texttt{R1: X X □ X □ (1,2,4) {-} I said 1,2,4. Correct.} \\
\texttt{R2: □ X X □ X (2,3,5) {-} I said 2,3,5. Correct.} \\
\texttt{R3: X □ □ X □ (1,4) {-} I said 1,4. Correct.} \\
\texttt{R4: □ X □ □ X (2,5) {-} I said 2,5. Correct.} \\
\texttt{R5: X □ X □ X (1,3,5) {-} I said 1,3,5. Correct.} \\
\\ 
\texttt{Wait, if my description was factually correct based on the provided text, why did I lose?} \\
\texttt{Ah, looking closely at the prompt's target grid again:} \\
\texttt{Target:} \\
\texttt{X X □ X □} \\
\texttt{□ X X □ X} \\
\texttt{X □ □ X □} \\
\texttt{□ X □ □ X} \\
\texttt{X □ X □ X} \\
\\ 
\texttt{My description:} \\
\texttt{R1: 1,2,4. (Matches X X . X .)} \\
\texttt{R2: 2,3,5. (Matches . X X . X)} \\
\texttt{R3: 1,4. (Matches X . . X .)} \\
\texttt{R4: 2,5. (Matches . X . . X)} \\
\texttt{R5: 1,3,5. (Matches X . X . X)} \\
\\ 
\texttt{Is it possible the failure was due to ambiguity or not distinguishing enough from distractors?} \\
\texttt{Let's check Distractor 1 R3: `X □ X X □` (1,3,4). My description "1 and 4" distinguishes it.} \\
\texttt{Distractor 2 R3: `X X □ □ X` (1,2,5). My description "1 and 4" distinguishes it.} \\
\\ 
\texttt{Perhaps the error was in the *format* or *conciseness*? The instruction says "Generate the referring expression starting with the tag 'Expression: ' ... Omit any other text." My output started with "Expression: ". That seems fine.} \\
\texttt{Maybe I misread the grid in the prompt history?} \\
\texttt{Let's look at the "LOSE" reason again: "failed to achieve the target state". In these tasks, usually, the model generates a string, and the evaluator checks if it uniquely} \\
\\ 
\\ 
\texttt{You are given three grids, where each of them is 5 by 5 in size.} \\
\texttt{Grids have empty cells marked with "▢" and filled cells marked with "X".} \\
\texttt{Your task is to generate a referring expression that best describes the target grid while distinguishing it from the two other distractor grids.} \\
\texttt{The first grid is the target grid, and the following two grids are the distractors.} \\
\\ 
\texttt{Target grid:} \\
\\ 
\texttt{X X □ X □} \\
\texttt{□ X X □ X} \\
\texttt{X □ □ X □} \\
\texttt{□ X □ □ X} \\
\texttt{X □ X □ X} \\
\\ 
\texttt{Distractor grid 1:} \\
\\ 
\texttt{X X □ □ X} \\
\texttt{□ X X □ X} \\
\texttt{X □ X X □} \\
\texttt{□ X □ X X} \\
\texttt{X □ □ X X} \\
\\ 
\texttt{Distractor grid 2:} \\
\\ 
\texttt{X □ X X □} \\
\texttt{□ X □ X X} \\
\texttt{X X □ □ X} \\
\texttt{□ X X □ X} \\
\texttt{X □ X □ □} \\
\\ 
\texttt{Instruction: Describe the target grid.} \\
\texttt{Generate the referring expression starting with the tag "Expression: " for the given target grid. Omit any other text.} \\
            }
        }
    }
    & \\ \\

    \theutterance \stepcounter{utterance}  
    & \multicolumn{4}{p{0.6\linewidth}}{
        \cellcolor[rgb]{0.9,0.9,0.9}{
            \makecell[{{p{\linewidth}}}]{
                \texttt{\tiny{[P1$\rangle$GM]}}
                \texttt{Expression: The grid where Row 1 has filled cells at columns 1, 2, and 4; Row 2 at columns 2, 3, and 5; Row 3 at columns 1 and 4; Row 4 at columns 2 and 5; and Row 5 at columns 1, 3, and 5.} \\
            }
        }
    }
    & & \\ \\

    \theutterance \stepcounter{utterance}  
    & & & \multicolumn{2}{p{0.3\linewidth}}{
        \cellcolor[rgb]{0.9,0.9,0.9}{
            \makecell[{{p{\linewidth}}}]{
                \texttt{\tiny{[GM$|$GM]}}
                \texttt{Expression: The grid where Row 1 has filled cells at columns 1, 2, and 4; Row 2 at columns 2, 3, and 5; Row 3 at columns 1 and 4; Row 4 at columns 2 and 5; and Row 5 at columns 1, 3, and 5.} \\
            }
        }
    }
    & & \\ \\

    \theutterance \stepcounter{utterance}  
    & & \multicolumn{4}{p{0.6\linewidth}}{
        \cellcolor[rgb]{0.9,0.9,0.9}{
            \makecell[{{p{\linewidth}}}]{
                \texttt{\tiny{[GM$\rangle$P2]}}
                \texttt{You have attempted this same task before and did not succeed. These are the plans you wrote after those attempts.} \\
\texttt{Trial \#1: New plan: I incorrectly identified the First grid as the match, but upon re{-}evaluating the expression against the First grid, I noticed a discrepancy in the third row (the expression requires filled cells at positions 1 and 4, but the First grid has them at 1 and 3). I failed to systematically verify every row of every grid against the specific positional constraints provided in the expression. Next time, I will methodically compare the expression's row{-}by{-}row requirements with each grid's layout, specifically checking the exact column indices for filled cells in the Third, Fourth, and Fifth rows to ensure a perfect match before selecting an answer. The correct grid is the Second one.} \\
\\ 
\\ 
\texttt{You are given three grids, where each of them is 5 by 5 in size.} \\
\texttt{Grids have empty cells marked with "▢" and filled cells marked with "X".} \\
\texttt{You are also given a referring expression that describes one of the given grids.} \\
\texttt{Your task is to select the grid that matches the given referring expression.} \\
\texttt{Generate only the number (in text) of the grid that matches the given expression by selecting first, second, or third.} \\
\\ 
\texttt{First:} \\
\\ 
\texttt{X X □ □ X} \\
\texttt{□ X X □ X} \\
\texttt{X □ X X □} \\
\texttt{□ X □ X X} \\
\texttt{X □ □ X X} \\
\\ 
\texttt{Second:} \\
\\ 
\texttt{X X □ X □} \\
\texttt{□ X X □ X} \\
\texttt{X □ □ X □} \\
\texttt{□ X □ □ X} \\
\texttt{X □ X □ X} \\
\\ 
\texttt{Third:} \\
\\ 
\texttt{X □ X X □} \\
\texttt{□ X □ X X} \\
\texttt{X X □ □ X} \\
\texttt{□ X X □ X} \\
\texttt{X □ X □ □} \\
\\ 
\texttt{Expression: The grid where Row 1 has filled cells at columns 1, 2, and 4; Row 2 at columns 2, 3, and 5; Row 3 at columns 1 and 4; Row 4 at columns 2 and 5; and Row 5 at columns 1, 3, and 5.} \\
\texttt{Question: Which grid does the expression refer to?} \\
\texttt{Start with the tag "Answer: ", followed by your selection. Omit any other text.} \\
            }
        }
    }
    & \\ \\

    \theutterance \stepcounter{utterance}  
    & & & \multicolumn{4}{p{0.6\linewidth}}{
        \cellcolor[rgb]{0.9,0.9,0.9}{
            \makecell[{{p{\linewidth}}}]{
                \texttt{\tiny{[GM$\langle$P2]}}
                \texttt{Answer: second} \\
            }
        }
    }
     \\ \\

    \theutterance \stepcounter{utterance}  
    & & & \multicolumn{2}{p{0.3\linewidth}}{
        \cellcolor[rgb]{0.9,0.9,0.9}{
            \makecell[{{p{\linewidth}}}]{
                \texttt{\tiny{[GM$|$GM]}}
                \texttt{second} \\
            }
        }
    }
    & & \\ \\

\end{supertabular}
}

\end{document}
